\section*{Ruta didáctica del manual}
\addcontentsline{toc}{section}{Ruta didáctica del manual}

Este manual no debe leerse como una colección de prácticas. Su propósito es acompañar al estudiante en una secuencia de decisiones experimentales: primero observar un fenómeno, después convertirlo en una pregunta, más tarde medirlo con cuidado y finalmente decidir con evidencia.

\begin{important}
La pregunta que guía el manual es simple: ¿cómo puede un ingeniero distinguir entre una impresión, una diferencia real y una conclusión defendible?
\end{important}

La lectura se organiza en cinco movimientos:

\begin{enumerate}
\item \textbf{Observar sin engañarse.} El estudiante comienza con fenómenos cotidianos y aprende que una sola observación no basta para afirmar una causa.
\item \textbf{Medir la variabilidad.} Antes de comparar tratamientos, se aprende a mirar muestras, medias, dispersión, datos faltantes y sesgos.
\item \textbf{Comparar tratamientos.} El ANOVA aparece como una forma de razonar con variabilidad, no como una tabla que se llena mecánicamente.
\item \textbf{Diseñar combinaciones.} Los diseños factoriales permiten estudiar factores, niveles, interacciones, blancos y réplicas.
\item \textbf{Decidir en sistemas reales.} Taguchi, modelos, ANCOVA, MANOVA y proyectos integradores se presentan como extensiones para problemas donde el experimento ideal no siempre es posible.
\end{enumerate}

Cada práctica debe cumplir una función dentro de esa ruta. Si una actividad no ayuda a formular mejor la pregunta, medir mejor, comparar mejor o decidir mejor, entonces debe reescribirse o moverse.

\begin{observacion}
Los documentos originales son materia prima. El estudiante no necesita saber de qué archivo provino una idea; necesita entender por qué aparece en ese momento del aprendizaje.
\end{observacion}

\newpage
